\section{Consumo de Servicios Web Externos y SOAP}
\subsection{SOAP vs. REST: análisis técnico comparativo}

SOAP (\textit{Simple Object Access Protocol}) es un protocolo de intercambio de mensajes basado en XML que describe
sus contratos mediante \textbf{WSDL} (\textit{Web Services Description Language}). Los mensajes viajan sobre HTTP,
SMTP o TCP y cada operación se identifica mediante una \textit{SOAPAction}. REST, en cambio, es un estilo
arquitectónico que usa cualquier formato de representación (JSON preferido), se apoya únicamente en HTTP y su
descripción es opcional (OpenAPI). En la Tabla~\ref{tab:soap} se comparan ambos enfoques.

\begin{table}[H]
\centering
\small
\setlength{\tabcolsep}{5pt}
\begin{tabular}{@{}p{3.2cm}p{4.6cm}p{4.6cm}@{}}
\toprule
\textbf{Criterio} & \textbf{SOAP} & \textbf{REST} \\
\midrule
Formato del mensaje & XML (Envelope) & JSON / XML / texto \\
Overhead de serialización & Alto (envoltorio e.g. envelopes) & Bajo \\
Tipado & Fuerte (esquemas XSD) & Débil (schemas opcionales) \\
Interoperabilidad & Muy alta (estándares WS-*) & Alta (más simple) \\
Facilidad desde JavaScript & Baja (requiere librerías) & Alta (fetch/Ajax nativo) \\
Casos de uso vigentes & Banca, sector público, SRI, SENESCYT & Aplicaciones web modernas \\
\bottomrule
\end{tabular}
\caption{Comparación SOAP vs. REST para el Sistema de Pre-Sustentaciones.}
\label{tab:soap}
\end{table}

En el ámbito universitario ecuatoriano, SOAP sigue vigente en sistemas gubernamentales como la verificación de
títulos en SENESCYT y facturación electrónica del SRI. Para el Sistema de Pre-Sustentaciones UTEQ, REST es la
elección adecuada para la comunicación entre el frontend Angular y el backend Spring Boot.

\subsection{Clientes HTTP del lado del servidor}

Para el consumo de APIs externas desde el servidor, los frameworks ofrecen clientes dedicados: \textbf{GuzzleHTTP}
en PHP, \textbf{HttpClient} en .NET, y \textbf{RestTemplate / WebClient} en Java/Spring. El \texttt{WebClient}
reactivo ofrece manejo de \textit{timeout}, reintentos con \textit{backoff exponencial} y concurrencia no
bloqueante. En la integración futura del PFC con una API externa (por ejemplo, consulta del calendario académico o
validación de docentes) se implementará \texttt{WebClient} con configuración de \textit{connect timeout} y
\textit{read timeout}, manejo de errores 4xx/5xx mediante \texttt{onErrorResume} y política de reintentos.

\subsection{Caché de respuestas externas}

Las APIs externas imponen límites de \textit{rate limiting}; por ello sus respuestas deben cachearse para evitar
superar cuotas y reducir la latencia. El patrón \textbf{cache-aside} consiste en consultar primero Redis; si hay
\textit{hit}, se devuelve la respuesta cacheada; si hay \textit{miss}, se consulta el origen, se almacena en
caché con un TTL y se responde. La configuración usará \texttt{@Cacheable("api_externa")} y un TTL por defecto en Redis.

\subsection{Tendencias: Jamstack y APIs en 2025}

El patrón \textbf{Jamstack} (JavaScript + APIs + Markup) separa frontend y backend. El Sistema de Pre-Sustentaciones
sigue una arquitectura SPA con Angular 21 y un backend de API Spring Boot 3.2.1, alineado con esta tendencia moderna.
